\chapter{Myomectomy}
\section*{Overview}
Myomectomy removes uterine fibroids while preserving the uterus.
\section*{Common Instruments and Equipment}
The exact setup varies by technique and setting. Common items may include:
\begin{itemize}[nosep,leftmargin=*]
  \item open, laparoscopic, or hysteroscopic gynecologic instruments
  \item uterine manipulators or hysteroscope depending on approach
  \item electrosurgical device and uterine closure sutures
\end{itemize}

\section*{Why It Is Done}
It may relieve heavy bleeding, pelvic pressure, pain, infertility-related distortion, or bulk symptoms when uterine preservation is desired.
\section*{How It Is Performed}
Fibroids may be removed hysteroscopically, laparoscopically, robotically, or through an abdominal incision depending on size, number, and location.
\section*{Contraindications}
No single contraindication applies to every patient; the treating team assesses whether the procedure is suitable for the individual and the reason for treatment.

\section*{Precautions}
The care team reviews relevant health conditions, medicines, allergies, and procedure-specific risks, and explains any changes needed for the individual.

\section*{Risks and Recovery}
Bleeding, infection, adhesions, recurrence, injury to nearby organs, and possible conversion to hysterectomy are possible. Future pregnancy planning requires specialist guidance.
\section*{References}
\begin{enumerate}\item MedlinePlus. Uterine Fibroids. U.S. National Library of Medicine; 2025.\item Current Medical Diagnosis and Treatment. McGraw-Hill; 2025.\end{enumerate}
\clearpage
